\documentclass[10pt,a4paper]{amsart}
\usepackage[margin=19mm]{geometry}
\usepackage{amsmath,amssymb,mathtools}
\usepackage[scr=rsfs,cal=euler]{mathalfa}
\usepackage{xfrac}
\usepackage[unicode,hidelinks,bookmarks=false]{hyperref}
\newcommand{\ie}{i.e.\ }
\newcommand{\cf}{cf.\ }
\newcommand{\resp}{respectively\ }
\newcommand{\Subsection}{\subsection}
\providecommand{\nobreakdash}{\nobreak\hbox{-}}
\setlength{\emergencystretch}{2em}
\multlinegap0pt
\usepackage{amssymb}
\usepackage[shortlabels]{enumitem}
\usepackage[normalem]{ulem}
\usepackage{tikz}
\newtheorem{lemma}{Lemma}
\usepackage{cleveref}
\crefname{lemma}{Lemma}{Lemmas}
\newcommand{\Ends}{\mathrm{Ends}}
\newcommand{\N}{\mathbb{N}}
\renewcommand{\l}{\lambda}
\DeclareMathOperator*{\limset}{\mathrm{lim}\; \mathrm{set}}
\newcommand{\mc}{\mathcal}
\newcommand{\orb}{\mathcal{O}}
\newcommand\ssm{\smallsetminus}
\newcommand{\vep}{\varepsilon}
\title{Cannon--Thurston maps for Anosov foliations}
\author{Ellis Buckminster}
\date{Source: arXiv:2604.21201v1 (CC BY 4.0)}
\begin{document}
\maketitle

\noindent\emph{Source and licence.} Cannon--Thurston maps for Anosov foliations. Version arXiv:2604.21201v1 (CC BY 4.0), distributed by arXiv under the Creative Commons Attribution 4.0 International licence, http://creativecommons.org/licenses/by/4.0/, as recorded in the arXiv licence metadata for this exact version and shown on the abstract page https://arxiv.org/abs/2604.21201v1. Full licence notice: \texttt{licenses/arxiv-2604.21201-CC-BY-4.0.txt}.\par\smallskip

\noindent\textit{Editorial scope.} Counted unit: the article's lemma lem:linf\_ray (source0.tex lines 1476--1481). The article's complete original proof of it is delivered with it (source0.tex lines 1483--1532, one \texttt{proof} environment of 50 lines). No mathematics is written, repaired or re-derived: the statement and the proof are the article's own lines, in the article's own notation, and no retained line is substituted or re-flowed.  No further source-local context is needed: every cross-reference of every retained line resolves inside the two delivered spans.  One declared adaptation: exactly 2 ancillary strings inside the retained spans are omitted (image-inclusion commands and, where the source repeats a label, the repeated label definitions) and no other line differs from the source; the figure environments, with their placement, captions and labels, are kept, and the package records the omitted strings in fidelity receipts bound to the affected bodies.  No retained line contains a citation, so the wrapper carries no bibliography and no bibliography entry.  The retained lines contain the article's own diagram code, which the wrapper typesets with the standard tikz packages; no image file is delivered and the package declares no asset dependency.  Editorial formatting only, in addition: an AMS \texttt{amsart} preamble that restates the article's own declarations and macros used by the retained lines, namely the environments \texttt{lemma} on separate counters, together with the standard packages those lines need.  The retained environments print in this document as Lemma 1 in order of appearance, whereas the article prints the same items with the numbering of its own full document.  The complete native source of the article as distributed in the arXiv e-print for this exact version is bundled unchanged as \texttt{sources/199005/source0.tex}.
\begin{lemma}\label{lem:linf_ray}
	Let $\vep\in\Ends(\mc L^u)$ such that $\vep\neq\l_\infty$. Let $\gamma^\vep_j$ be the zigzag ray from $\l'_j$ to $\vep$. Then 
	\[
	\l_\infty\in\limset_{k\to\infty}\gamma^\vep_k.
	\]
\end{lemma}

\begin{proof}
	Let $U$ be a neighborhood of $\l_\infty$. We need to show that there exists some $K_U\in\N$ such that for all $k\geq K_U$, $\gamma^\vep_k\cap U\neq\varnothing$. The first thing to notice is that for all $k$,
	\[
	\gamma_k^\vep=\l'_k\cup\{\l^u\in\mc L^u\;\mid\;\ell^u\text{ separates }\ell'_k\text{ and }f(\vep)\text{ in }\overline{\orb}\}.
	\]
	
	We split the proof into the following three cases, either 
	\begin{enumerate}
		\item $z_\infty\in\orb$,
		\item $z_\infty\in\partial\orb^u$, or
		\item $z_\infty\in\partial\orb\ssm\partial\orb^u$.
	\end{enumerate}
	
	\begin{figure}[h]
		\centering
		
		\caption{On the left, case (1) in the proof of \Cref{lem:linf_ray}, with $V$ shaded in yellow. On the right is case (3), with $H_{N+2}$ shaded in yellow.}\label{fig:linf_ray}
	\end{figure}

	In case (1), let $V$ be an open foliation chart around $z_\infty$ such that the projection of $V$ to $\mc L^u$ is contained in $U$. Since 
	\[
	\lim_{k\to\infty}z_k'=z_\infty,
	\]
	past some $K_V\in\N$ we have $z_k'\in V$, so $\ell'_k$ intersects $V$. Then using a stable arc $t^s$ in $V$, we may shift $\ell'_k$ slightly to one side to produce an unstable leaf $t^u$ separating $\ell'_k$ from $f(\vep)$ and intersecting $V$. See \Cref{fig:linf_ray}. Thus, $K_U=K_V$ works.
	
	For case (2), let $\ell_1^u,\ell^s_1,\ldots,\ell^u_n,\ell^s_n$ be the leaves of $\orb^{s/u}$ ending at the point $z_\infty$, and arranged in order as in \Cref{fig:linf_ray2} (it might be the case that $\ell_1^u$ or $\ell^s_n$ does not exist). A neighborhood $V$ for $z_\infty$ in $\overline{\orb}$ is given by the set bound by an arc of $\partial\orb$, together with segments of leaves $t^s_1,t^u_s,\ldots,t^s_n,t^u_n$ as in \Cref{fig:linf_ray2}. We can shrink this neighborhood to a neighborhood $V_U$ such that the projection of $V_U$ to $\mc L^u$ is contained in $U$ by shrinking the $t^{s/u}_j$'s along the $\ell^{s/u}_j$'s and by removing halfspaces bound by unstable leaves. We further shrink $V_U$ so that $f(\vep)\notin V_U$.

	
	Since 
	\[
	\lim_{k\to\infty}z_k'=z_\infty,
	\]
	we must have that past some $K_V\in\N$, $z_k'\in V_U$.  Either 
	\[
	z_k'\in\overline{\mc S^u(t_1^s\cup\ldots\cup t^s_n)},
	\]
	or $z_k'$ is separated from $f(\vep)$ by a leaf $\ell^u$ nonseparated from a leaf in $\mc S^u(t_1^s\cup\ldots\cup t^s_n)$. In the first case, we again can slightly nudge the unstable leaf $z_k'$ is on along the appropriate $t^s_j$ and toward $f(\vep)$ to produce a leaf $\ell^u$ in $\gamma_k^\vep\cap U$. In the second case, $\ell^u\in\gamma_k^\vep\cap U$.
	
		\begin{figure}[h]
		\centering
		
		\caption{Case (2) for the proof of \Cref{lem:linf_ray}. The set $V_U$ is shaded in yellow. Two options for $z_k$ and the corresponding choices of $\ell^u$ are shown.}\label{fig:linf_ray2}
	\end{figure}
	
	For case (3), a neighborhood basis for $z_\infty$ in $\overline{\orb}$ is given by the half-planes $H_n$ bound by a set of leaves 
	\[
	\{\ell^u_n\;\mid\;n\in\N,\ell^u_n\in\orb^u\},
	\] as in \Cref{fig:linf_ray}. Let $V$ be one of these halfspaces. As in case (2), we shrink $V$ to an open set $V_U$ whose projection to $\mc L^u$ is contained in $U$ by picking a smaller halfspace $H_N$ in the above set and removing some halfspaces bound by unstable leaves. As in the previous cases, the $z_k'$'s eventually lie in $H_{N+2}$, so $\ell^u_{N+1}$ is a leaf in $\gamma_k^\vep\cap U$.

\end{proof}

\end{document}
